\chapter{Discussion and Conclusions}\label{ch:conclusions}

\section{Answers to the research questions}\label{sec:conclusions-findings}

\paragraph{RQ1: the joint score and its couplings.} For the stationary AR(1)
chain under coordinatewise OU corruption the joint score is linear and known in
closed form, $\score(x,t) = -\Qt x$ with
$\Qt = (e^{-2t}\Sigz + \Deltat I)^{-1}$, equal by the score--posterior identity
to $(e^{-t}\E[a|x] - x)/\Deltat$. Its coupling structure has a closed-form
lifecycle: exactly tridiagonal at $t = 0$, filling one further diagonal per
order in $t$ by the band-fill law \eqref{eq:g-bandfill}, maximally coupled at
intermediate times, and returning to the identity at rate $e^{-2t}$. In the
bulk a single quantity $q(\alpha,t) \in (0, |\alpha|]$ governs the posterior
correlation length. This is exact for the Gaussian AR(1) model and for no
other: it is the linearity of that model, not the chain structure alone, that
makes the score a matrix.

\paragraph{RQ2: belief propagation and the Kalman smoother.} The posterior of
the noised chain lives on a tree-shaped factor graph
(Figure~\ref{fig:posterior-factor-graph}), because the channel corrupts each
coordinate independently and every observation therefore enters through a unary
factor. One forward and one backward sweep compute all smoothing marginals. On
the Gaussian chain the messages close algebraically in the Gaussian family, and
the recursion is the Kalman filter followed by the RTS smoother, update by
update. The two routes agree to $10^{-15}$--$10^{-14}$ across the tested
$(\nsites, \alpha, t)$ grid, with an independently written implementation
reproducing the agreement. Tree exactness is a property of the graph and holds
for any innovation law; closure in a finite-dimensional family is a property of
the model, and holds here only because the model is Gaussian.

\paragraph{RQ3: what Gaussian messages cost.} On Laplace
chains, Gaussian-projected BP measured against validated grid quadrature has
median relative score error $0.24$ at $t = 0.05$, falling below $10^{-2}$ by
$t \approx 1$ and below $10^{-3}$ by $t \approx 2.1$; no trial exceeds $10\%$
relative error beyond $t \approx 0.98$, and the score direction is preserved far
better than its magnitude (median cosine $\geq 0.976$ throughout, worst single
trial $0.879$). What that number measures is identified exactly by
Theorem~\ref{thm:l-lmmse}: moment-projected sum--product returns the LMMSE
estimator of the covariance-matched Gaussian model, so the comparison is full
non-Gaussian posterior inference against the best second-order linear
estimator, and nothing else. The error is amplified at small $t$ by the exact
factor $e^{-t}/\Deltat$, is worst for bimodal innovation laws, and decreases
with stronger chain correlation. The reference is numerical, so these are
measurements within its validated resolution regime rather than formulas.

\paragraph{RQ4: how fast remote evidence stops mattering.} The influence of
evidence at distance $d$ on the posterior mean decays exactly as
$q(\alpha,t)^d$. Separately, the error of the radius-$r$ windowed score
estimator is \emph{measured} to decay at the same geometric rate over
$r = 2,\dots,13$. The two are distinct claims: windowing changes the boundary
conditions of the recursion in a way the influence-coefficient argument does not
cover, so the second is a numerical law over the tested range and not a
corollary of the first. At equal locality budget, truncating the inference is
more accurate than truncating the score matrix at small $t$ ($12.3\%$ against
$21.0\%$ at $t = 0.05$). On the same model the per-node AMP/TAP closure returns
the exact score --- the mean fixed point is closure-independent --- but a
different posterior variance, with existence boundary, breakdown time
$t_c(\alpha)$ and critical coupling $\alpha_c = \sqrt{2}-1$ in closed form
(Appendix~\ref{app:amp}).\label{sec:conclusions-amp}

\paragraph{RQ5: recovering an unknown transition, and what must be checked
first.} Fisher's identity turns one complete forward--backward pass into the
exact gradient of the marginal log-likelihood, with no differentiation through
the recursion, so the E-step is structurally exact for any innovation law; the
implementation adds controlled grid quadrature and a generalised M-step. The
checks that recovery requires are the substance of the answer rather than a
caveat on it. The channel acts twice --- on the local EM rate through the
fraction of missing information, and on the observed Fisher information at
finite sample size --- and it acts unevenly across coordinates: the fitted
innovation shape settles at a median of $\convshapemed$ updates against
$\convrhomed$ for the correlation, so a stopping rule read off the correlation
trace certifies a kernel that has not converged in the coordinate every
non-Gaussian claim depends on. A fixed iteration budget therefore compares
convergence rates, not estimators, and any claim resting on fitted shape at a
fixed budget must say so.

\paragraph{RQ6: what supplying the Markov assumption is worth.} At equal
training-set size the structured estimator attains $\ratiolo$--$\ratiohi\times$
lower pointwise denoising error than networks trained by denoising score
matching on identical data, across the \nsizesused{} sizes tested. Against a
baseline that is given locality and weight sharing --- the window head selected
by a screen over three architectures on disjoint seeds --- the advantage is
$\structratiolo$--$\structratiohi\times$ (Figure~\ref{fig:em-value-of-structure}).
Three scope conditions travel with those numbers. The two arms are asymmetric in
supplied information by construction, so this measures what a correct structural
assumption is worth on a family where it holds exactly, and is not a ranking of
learning algorithms. The difference between the two ratios is not a
decomposition into an architecture term and an estimation term: architecture,
capacity, parameter sharing, optimisation and implicit regularisation change
together when the baseline is swapped. And both are ratios of errors at fixed
data, not ratios of the data each method needs to reach a fixed error --- the
curves do not cross in the measured range, so no data-equivalence factor
follows.

\section{What is new}

\paragraph{The exact structure of a chain diffusion score.} The joint score of a
noised Markov chain is characterised completely for the Gaussian case: its
closed form, its eigenbasis at every diffusion time, the law by which its
couplings fill and then decay, and the single parameter that controls posterior
influence range. The identification of chain BP with the Kalman--RTS recursion
makes the two classical literatures one object rather than two analogous ones.

\paragraph{An exact reading of what Gaussian closure computes.} Where previous
practice measured the error of moment-projected message passing as the error of
an approximation, Theorem~\ref{thm:l-lmmse} identifies the estimator it returns.
The measured Laplace gap is therefore the price of discarding everything beyond
second moments --- a property of the model class, shared by every innovation law
with the same first two moments --- and not a property of one implementation.
The validated grid reference that makes the measurement possible is itself a
contribution, since the functional recursion has no closed form.

\paragraph{Transition learning, and the measured value of structure.} The
transition law is estimated from noised sequences alone by exact chain inference
inside EM, and the resulting estimator is compared against trained networks at
equal data under a protocol that selects both arms on held-out validation. Two
boundaries of the advantage are mapped rather than asserted: mixture capacity
beyond a single Gaussian innovation buys nothing for this estimator under the
grid and cap implemented here, and a chain-structured estimator partially
accommodates a rank-one violation of the Markov assumption while failing on a
long-range one at a coupling strength of roughly one tenth. The convergence
measurement that separates rate from estimator is a methodological result in its
own right, and it is what withdrew an earlier reading of the generative
comparison.

\section{Limitations}

\begin{enumerate}
  \item \textbf{Model class.} The exact solutions concern a one-dimensional
  state per frame and linear AR(1) dynamics, extended to a two-dimensional
  nonlinear state only in Chapter~\ref{ch:ring}. The locality laws are proved
  for the Gaussian bulk; their non-Gaussian counterparts are only structurally
  delimited, and no low-rank structure is asserted for the Laplace
  $\nsites = 2$ posterior Hessian beyond the narrow regime where it holds.

  \item \textbf{The numerical reference.} The non-Gaussian results are
  numerical. Grid BP is a reference only within its validated resolution regime
  ($\sqrt{2t} \gtrsim 3h$ on the grid spacing $h$), and the Laplace closure
  error is a measured curve rather than a formula.

  \item \textbf{Optimisation and finite seeds.} The learning results are
  statistical rather than exact: estimates over a finite number of seeds,
  reported with paired uncertainty, several resolved at only one of the two
  training sizes tested. Every capacity cell above $C = \capcappedfrom$ stops at
  an iteration cap rather than at tolerance, so a statistical disadvantage of
  larger mixtures is not separated from an optimisation or a discretisation one.
  The misspecification study probes exactly two violations of the Markov
  assumption; it establishes that the two are not interchangeable in effect, not
  a general theory of how much non-Markov structure a chain estimator tolerates.

  \item \textbf{Generative fidelity, and transfer.} Whether pointwise accuracy
  and generative fidelity can move in opposite directions within one estimator
  is unresolved here: the comparison that addresses it evaluates procedures
  whose fitted shape coordinates are unequally settled, on the single coordinate
  the generative metric scores, so its reading is withdrawn rather than
  reversed. More broadly, these are properties of one controlled model. Whether
  trained score networks on real sequence data inherit them is untested here.
  The image and video extensions explored in the underlying research package are
  outside the scope of this thesis and are not reported.
\end{enumerate}

\section{Future work}

Four directions, in the order in which they would most change what this thesis
can claim.

\emph{First, the reverse-generation comparison, done with converged fits.}
Section~\ref{sec:em-generative} integrates the reverse SDE under the estimated
EM--BP score and under two trained neural scores, but on fits the convergence
analysis shows to be unequally settled. Selecting both arms on a convergence
criterion rather than a fixed iteration count, then comparing, would answer the
question the experiment was built for and would close the one question this
thesis leaves open rather than answers.

\emph{Second, a grid-by-capacity design.} The capacity result leaves two
mechanisms entangled. A finer grid at fixed capacity would say how much of the
larger mixtures' disadvantage is their fits falling below the grid's resolution
rather than their overfitting the data, since only the second should survive a
mesh refinement; a budget ladder at fixed capacity would do the same for the
iteration cap.

\emph{Third, richer departures from the Markov assumption.} The present study
probes two violations at fixed strengths. A paired design across violation
strengths, and further violation types beyond a global latent and long-range
coupling, would test whether the finding that the two mechanisms are not
interchangeable generalises. Discrete chains are untouched here: with a finite
alphabet the BP messages are finite vectors and inference is exact with no
closure at all, which would give a second solvable benchmark alongside the
Gaussian one.

\emph{Fourth, extension to image and video models.} The reverse-dynamics
comparison used a small chain length and one innovation family. Whether the
structural results survive longer sequences, other non-Gaussian innovation laws,
and the move to data with spatial as well as temporal structure is the question
that decides how far any of this carries beyond the solvable setting.
